\appendix
\setcounter{table}{0}
\renewcommand{\thetable}{A\arabic{table}}
\setcounter{figure}{0}
\renewcommand{\thefigure}{A\arabic{figure}}
% \noindent\textbf{Annotators.}\qquad
% All annotators are students from high level Chinese universities and we have tested their language skills to ensure that they can produce high quality annotations. During the annotation process, some of the annotators showed extremely outstanding linguistic skills, which also provided us with selection criteria for segmenting the dataset.\\
% \noindent\textbf{Pre-processing.}\qquad
% Before annotation, we manually cleaned the collected videos with two stages of clips. In the first stage, we cut all videos into content-independent clips, for example, we cut KT into individual performances. In the second stage, we first remove all videos that are unethical and contain sensitive information we cleaned the Humor videos to remove clips that do not meet ethical standards and have sensitive information, we do not operate on KT clips, only the score of each performance clip is recorded, and we subdivided the magic videos into clips to ensure that each clip includes only 1-3 coherent magic effects to avoid exceeding the model's understanding ability.\\
% \noindent\textbf{Tasks Annotation.}\qquad
% 1) We have developed detailed annotation documentation, and all annotators undergo rigorous training before performing the actual annotation. 
% 2) For the tasks in free-text, in Humor QA and MagicQA we suggested answers between 2-4 short sentences, in CreativeQA we would like a longer answer from the annotators. We asked the annotators to use concise and objective statements and avoid subjective adverbs and adjectives to ensure that the answers were in a similar language style. 
% 3) Annotators can question the counter-intuitive properties of video content or the generalizability of answers during the annotation process, and we will document such issues and address them in the next stage.
% 4) Cross-validation: We arranged another group of people who did not participate in the labeling to cross-validate the QA pair into three grades: completely wrong, ambiguous or too subjective, and no problem, and the statistical results can be seen in Figure \ref{}. because people's subjective understanding of humor, creativity and magic are bound to differ, a small number of disagreements is a normal result, we fixed the completely wrong QA pair and kept the rest.
% \noindent\textbf{Post-Processing.}\qquad
% After completing a round of annotation, we will check and filter the annotations. To ensure the quality of annotations, annotations with low quality or ambiguity will be returned to the task annotation stage for the next round of annotation. After all task annotations are completed with high quality, we will enter the data augmentation phase.


\section{More Details of FunQA Dataset} \label{appendix_a: FunQA more details}
\subsection{Dataset Construction Pipeline} \label{appendix_a:dataset con}
\noindent\textbf{Video Selection}\quad
In constructing the dataset, we adhered to three principles to address the challenges in video understanding capabilities: our dataset, FunQA, is \textbf{visual centered} and emphasizes \textbf{counter-intuitive reasoning}, \textbf{spatial-temporal reasoning}. Based on these principles, we collect 4365 videos from 3 different art genres and created three subsets: HumorQA, CreativeQA, and MagicQA.\\
\noindent\textbf{HumorQA}\quad
HumorQA composed of 1,769 meticulously curated web videos, serves as a unique source of insight into human humor comprehension. Notably, it contains the shortest average video length of 7s among the three subsets. We believe that the human process of understanding humor is complex and deep, requiring a holistic understanding of the video and adding a degree of common sense to it. Psychological research has demonstrated that humor arises from the incongruity \cite{kant1987critique,sep-humor} between reality and expectations, flourishing with the skillful juxtaposition and transformation of events \cite{latta1999basic,boyd2004laughter,koestler2020act}. This makes humorous videos a valuable asset for the VideoQA dataset, anticipated to enhance a model's proficiency in integrating information and performing deep reasoning.\\
\noindent\textbf{CreativeQA}\quad 
CreativeQA is a collection of 927 videos averaging 48s in length from a TV show called Kasou Taishou \cite{ntv}. This program, showcasing original and novel skits performed by various amateur groups and judged by a panel, boasts a strong creative flair 
 \cite{runco2012standard}. The essence of the show lies in using a mix of people and props to mimic reality, with audiences deriving pleasure from information integration and comparison. We anticipate that the imitation nature of the show will challenge the model's capacity for information extraction, while the longer video length and need for understanding creativity will put to test the model's comprehension of spatial-temporal information.\\
\noindent\textbf{MagicQA}\quad
MagicQA encapsulates 1672 magic performance videos sourced from across the web, spanning various genres like camera magic, close-up magic, and stage magic. The essence of magic revolves around the creation of seemingly impossible illusions \cite{nelms2012magic}, employing diverse effects such as disappearance, creation, and transformation. These illusions are infused with abundant spatial-temporal information. Through this dataset, we aim to empower the model to not only track the ensuing changes in objects but also unravel the underlying mechanics \cite{lamont2005magic} of these transformations.

\begin{table*}[!t]
\centering
\caption{\textbf{Comparison between FunQA and other existing benchmarks (Complete Version).}
Compared to other datasets, FunQA revolves around the captivating realm of interesting and counter-intuitive videos. The tasks within FunQA are specifically designed to challenge the vision capabilities of models, requiring strong skills in producing an in-depth description, interpretation, and spatial-temporal reasoning.
Here we clarify the abbreviation in the table. For annotation type:  \includegraphics[width=0.02\textwidth]{figures/human.png} denotes Manual Annotation and  \includegraphics[width=0.02\textwidth]{figures/robot.png} for Automatic Annotation; \textbf{Avg Len} denotes video average length; \textbf{\# Clips} means number of video clips; \textbf{VC} for visual-centric, \textbf{Des.} for Description, \textbf{Exp.} for Explanation, \textbf{STR} for Spatial-temporal Reasoning, \textbf{MC} means Multiple Choice QA, and \textbf{OE} shows Open Ended QA with \textbf{Average Word Count} per response. 
}
\vspace{-5pt}
\label{T:whole_main}
\setlength\tabcolsep{10pt}
\resizebox{1.0\linewidth}{!}{
\begin{tabular}{llccccccccccl}
\toprule
\multirow{2}{*}{Dataset} & \multirow{2}{*}{Domain} & \multirow{2}{*}{
\includegraphics[width=0.02\textwidth]{figures/human.png} or \includegraphics[width=0.02\textwidth]{figures/robot.png}
} & \multicolumn{2}{c}{Video} & \multicolumn{7}{c}{Question Answer}  \\
\cmidrule(l){4-5} \cmidrule(l){6-12} 
 &  &  & Avg Len & \# Clips                      & \# QA  & VC & Des. & Exp. & STR & MC & OE   \\
 \toprule
TGIF-QA \cite{jang-IJCV-2019} 
& Social Media & \includegraphics[width=0.02\textwidth]{figures/robot.png} & 3s & 72K & 165K & \textcolor{darkgreen}{\ding{51}} & \textcolor{darkgreen}{\ding{51}} & \textcolor{darkred}{\ding{55}} & \textcolor{darkgreen}{\ding{51}} & \textcolor{darkred}{\ding{55}} & 2.1 \\

MSRVTT-QA \cite{xu2017video}
& Social Media & \includegraphics[width=0.02\textwidth]{figures/robot.png} & 15s & 10K & 244K & \textcolor{darkgreen}{\ding{51}} & \textcolor{darkred}{\ding{55}} & \textcolor{darkred}{\ding{55}} & \textcolor{darkgreen}{\ding{51}} & \textcolor{darkred}{\ding{55}} & 1.0 \\

ActivityNet-QA \cite{yu2019activitynetqa}
& Social Media & \includegraphics[width=0.02\textwidth]{figures/human.png} & 180s & 6K & 58K & \textcolor{darkgreen}{\ding{51}} & \textcolor{darkred}{\ding{55}} & \textcolor{darkred}{\ding{55}} & \textcolor{darkgreen}{\ding{51}} & \textcolor{darkred}{\ding{55}} & 1.9 \\

MSVD-QA \cite{xu2017video}
& Social Media & \includegraphics[width=0.02\textwidth]{figures/robot.png} & 10s & 2K & 51K & \textcolor{darkgreen}{\ding{51}} & \textcolor{darkred}{\ding{55}} & \textcolor{darkred}{\ding{55}} & \textcolor{darkred}{\ding{55}} & \textcolor{darkred}{\ding{55}} & 1.0\\

YouTube2Text-QA \cite{Ye_2017}
& Social Media & \includegraphics[width=0.02\textwidth]{figures/robot.png} & 10s & 10K & 123K & \textcolor{darkgreen}{\ding{51}} & \textcolor{darkred}{\ding{55}} & \textcolor{darkred}{\ding{55}} & \textcolor{darkgreen}{\ding{51}} & \textcolor{darkgreen}{\ding{51}} & N/A \\

AGQA \cite{grundemclaughlin2021agqa}                              
& Social Media & \includegraphics[width=0.02\textwidth]{figures/robot.png} & 30K & 10s & 192K & \textcolor{darkgreen}{\ding{51}} & \textcolor{darkred}{\ding{55}} & \textcolor{darkgreen}{\ding{51}} & \textcolor{darkgreen}{\ding{51}} & \textcolor{darkred}{\ding{55}} & TBD \\

AVQA \cite{yang2022avqa}                                          
& Social Media & \includegraphics[width=0.02\textwidth]{figures/human.png} & 60s & 9K & 57K & \textcolor{darkgreen}{\ding{51}} & \textcolor{darkred}{\ding{55}} & \textcolor{darkgreen}{\ding{51}} & \textcolor{darkgreen}{\ding{51}} & \textcolor{darkgreen}{\ding{51}} & N/A \\

\toprule

NExT-QA \cite{xiao2021nextqanext}            
& Daily life & \includegraphics[width=0.02\textwidth]{figures/human.png} & 44s & 5K & 52K & \textcolor{darkgreen}{\ding{51}} & \textcolor{darkgreen}{\ding{51}} & \textcolor{darkgreen}{\ding{51}} & \textcolor{darkgreen}{\ding{51}} & \textcolor{darkgreen}{\ding{51}} & 2.6 \\

Social-IQ \cite{zadeh2019social}                        
& Daily life & \includegraphics[width=0.02\textwidth]{figures/human.png} & 99s & 1K & 8K & \textcolor{darkgreen}{\ding{51}} & \textcolor{darkred}{\ding{55}} & \textcolor{darkgreen}{\ding{51}} & \textcolor{darkred}{\ding{55}} & \textcolor{darkgreen}{\ding{51}} & N/A \\

STAR \cite{wu2021star}                                            
& Daily life & \includegraphics[width=0.02\textwidth]{figures/robot.png} & - & 23K & 60K & \textcolor{darkgreen}{\ding{51}} & \textcolor{darkred}{\ding{55}} & \textcolor{darkred}{\ding{55}} & \textcolor{darkgreen}{\ding{51}} & \textcolor{darkgreen}{\ding{51}} & N/A  \\

FIBER \cite{castro2022fiber}                                      
& Daily life & \includegraphics[width=0.02\textwidth]{figures/human.png} & 10s & 28K & 2K & \textcolor{darkred}{\ding{55}} & \textcolor{darkgreen}{\ding{51}} & \textcolor{darkred}{\ding{55}} & \textcolor{darkgreen}{\ding{51}} & \textcolor{darkred}{\ding{55}} & TBD \\

\toprule

MovieQA \cite{tapaswi2016movieqa}                       & TV shows & \includegraphics[width=0.02\textwidth]{figures/robot.png} & 203s & 7K & 6K & \textcolor{darkred}{\ding{55}} & \textcolor{darkred}{\ding{55}} & \textcolor{darkgreen}{\ding{51}} & \textcolor{darkgreen}{\ding{51}} & \textcolor{darkgreen}{\ding{51}} & N/A \\

TVQA \cite{lei2018tvqa}    
& TV shows & \includegraphics[width=0.02\textwidth]{figures/human.png} & 76s & 22K & 153K & \textcolor{darkred}{\ding{55}} & \textcolor{darkred}{\ding{55}} & \textcolor{darkgreen}{\ding{51}} & \textcolor{darkred}{\ding{55}} & \textcolor{darkgreen}{\ding{51}} & N/A \\

TVQA+ \cite{lei2019tvqa}  
& TV shows & \includegraphics[width=0.02\textwidth]{figures/human.png} & 8s & 4K & 30K & \textcolor{darkred}{\ding{55}} & \textcolor{darkred}{\ding{55}} & \textcolor{darkgreen}{\ding{51}} & \textcolor{darkgreen}{\ding{51}} & \textcolor{darkgreen}{\ding{51}} & N/A \\

KnowIT-VQA \cite{garcia2020knowit}                                 
& TV shows & \includegraphics[width=0.02\textwidth]{figures/human.png} & 60s & 12K & 24K & \textcolor{darkgreen}{\ding{51}} & \textcolor{darkred}{\ding{55}} & \textcolor{darkgreen}{\ding{51}} & \textcolor{darkgreen}{\ding{51}} & \textcolor{darkgreen}{\ding{51}} & N/A \\

\toprule


SUTD-TrafficQA \cite{xu2021sutd}
& Traffic & \includegraphics[width=0.02\textwidth]{figures/human.png} & 5s & 10K & 623K & \textcolor{darkgreen}{\ding{51}} & \textcolor{darkred}{\ding{55}} & \textcolor{darkred}{\ding{55}} & \textcolor{darkgreen}{\ding{51}} & \textcolor{darkgreen}{\ding{51}} & N/A \\
MarioQA \cite{mun2017marioqa} 
& Games  & \includegraphics[width=0.02\textwidth]{figures/human.png} & 5s & 188K & 188K & \textcolor{darkgreen}{\ding{51}} & \textcolor{darkred}{\ding{55}} & \textcolor{darkgreen}{\ding{51}} & \textcolor{darkgreen}{\ding{51}} & \textcolor{darkred}{\ding{55}} & 2.0 \\

CLEVRER \cite{yi2020clevrer}     
& Synthetic Videos & \includegraphics[width=0.02\textwidth]{figures/human.png} & 5s & 20K & 305K & \textcolor{darkgreen}{\ding{51}} & \textcolor{darkred}{\ding{55}} & \textcolor{darkgreen}{\ding{51}} & \textcolor{darkgreen}{\ding{51}} & \textcolor{darkgreen}{\ding{51}} & N/A \\

Env-QA \cite{Gao_2021_ICCV}                                       
& Egocentric & \includegraphics[width=0.02\textwidth]{figures/human.png} & 20s & 23K & 85K & \textcolor{darkgreen}{\ding{51}} & \textcolor{darkred}{\ding{55}} & \textcolor{darkred}{\ding{55}} & \textcolor{darkgreen}{\ding{51}} & \textcolor{darkgreen}{\ding{51}} & N/A \\














% \midrule

\midrule
\textbf{HumorQA (Ours)     }                                     
& Daily life & \includegraphics[width=0.02\textwidth]{figures/human.png} & 7s & 2K & 141K & \textcolor{darkgreen}{\ding{51}} & \textcolor{darkgreen}{\ding{51}} & \textcolor{darkgreen}{\ding{51}} & \textcolor{darkgreen}{\ding{51}} & \textcolor{darkgreen}{\ding{51}} & \textbf{28.2}\\

\textbf{CreativeQA (Ours)   }                               
& Performance & \includegraphics[width=0.02\textwidth]{figures/human.png} & 48s & 1K & 79K & \textcolor{darkgreen}{\ding{51}} & \textcolor{darkgreen}{\ding{51}} & \textcolor{darkgreen}{\ding{51}} & \textcolor{darkgreen}{\ding{51}} & \textcolor{darkgreen}{\ding{51}} & \textbf{59.1}\\

\textbf{MagicQA (Ours)  }                              
& Magic shows & \includegraphics[width=0.02\textwidth]{figures/human.png} & 10s & 2K & 92K & \textcolor{darkgreen}{\ding{51}} & \textcolor{darkgreen}{\ding{51}} & \textcolor{darkgreen}{\ding{51}} & \textcolor{darkgreen}{\ding{51}} & \textcolor{darkgreen}{\ding{51}} & \textbf{27.6}\\

\textbf{FunQA (Ours)} 
& \textbf{Surprising Videos} 
& \includegraphics[width=0.02\textwidth]{figures/human.png} & 19s & 4K & \textbf{312K} & \textcolor{darkgreen}{\ding{51}} & \textcolor{darkgreen}{\ding{51}} & \textcolor{darkgreen}{\ding{51}} & \textcolor{darkgreen}{\ding{51}} & \textcolor{darkgreen}{\ding{51}} & \textbf{34.2} \\ 
\bottomrule 

\end{tabular}}
% \vspace{-0.5cm}
\end{table*}

\begin{figure}[h]
    \centering
    \includegraphics[width=\linewidth]{figures_appendix/M4.pdf}
    \vspace{-0.6cm}
    \caption{\textbf{Example of Magic Method Task.} During the annotation process, we discovered that even as humans, it is difficult to fully understand the complete principles behind the implementation of magic tricks right from the beginning.}
    \vspace{-0.4cm}
    \label{fig:M4}
\end{figure}

\begin{figure}[h]
    \centering
    \includegraphics[width=\linewidth]{figures_appendix/annotation_priciple.pdf}
    \vspace{-0.6cm}
    \caption{\textbf{FunQA Video Selection and Clip Principle.} We have a zero-tolerance policy regarding the inclusion of offensive content in our dataset. During the video sourcing process (video selection and video clipping stages), we ensure that such content is completely eliminated.}
    \vspace{-0.4cm}
    \label{fig:annotation-principle}
\end{figure}

\begin{figure*}[h]
    \centering
    \includegraphics[width=\linewidth]{figures_appendix/anno-error.pdf}
    \vspace{-0.6cm}
    \caption{\textbf{Examples of incorrect annotations and the suggested modifications.} Here we clarify the meanings of different cell colors. \textbf{Red}: Errors made by the annotator during the annotation process. \textbf{Yellow}: Types of annotation errors. \textbf{Green}: Reasons for the errors and suggested corrections.}
    \vspace{-0.4cm}
    \label{fig:annotation-error}
\end{figure*}

\begin{figure}[ht]
    \centering
    \includegraphics[width=\linewidth]{figures_appendix/prompt-trans.pdf}
    \vspace{-0.6cm}
    \caption{\textbf{GPT prompt for Translation and augmentation.} Under our carefully designed prompt, GPT-3.5 can automatically generate English translations and five synonymous sentences from our high-quality Chinese annotations, thereby expanding FunQA dataset.}
    \vspace{-0.4cm}
    \label{fig:prompt-trans}
\end{figure}
\begin{figure}[h]
    \centering
    \includegraphics[width=\linewidth]{figures_appendix/tran_err.pdf}
    \vspace{-0.6cm}
    \caption{\textbf{Issues in translation caused by Chinese-English cultural differences.} In this example of translating and expanding a humorous video title, GPT-3.5 failed to understand the meaning of the original Chinese title, and we filtered out such data and made manual changes.}
    \vspace{-0.4cm}
    \label{fig:trans-err}
\end{figure}


\noindent\textbf{Pre-processing \& Qualification}\quad
For videos related to humor and magic, we downloaded them from different streaming platforms, mostly in the form of compilations. For Creative videos, we downloaded 26 episodes publicly available from Tokyo TV in Japan. We provided rigorous training to the annotators to ensure high-quality video clips in the final compilation. Annotators who successfully completed the Clip task according to the requirements are considered qualified and can proceed to the next stage of annotation.

\noindent\textbf{Training \& Annotation}\quad
We conducted systematic training for all annotators who passed the previous round of annotation, focusing on different tasks: \\
For the \textit{Counter-intuitive Timestamp Localization Task (H1, C1, and M1)}, the annotation format is a pair of numbers enclosed in square brackets, [xxxx, xxxx]. We asked the annotators to record the time intervals in which they felt pleasure (or amusement or shock) while watching the video. \\
For the \textit{Task Detailed Description Task (H2, C2, and M2)}, we requested objective descriptions of what happened at [xxxx, xxxx], emphasizing a "what you see is what you get" approach. It is important to note that the annotations should only cover the selected time intervals and should not include subjective adverbs (such as vividly, vividly, or wildly). When describing characters or objects, be concise and add modifiers if there is ambiguity. An example of a poor label is "a man and a woman," while a relatively better example is "a man wearing a red hat and a woman wearing an apron." 

For the \textit{Counter-intuitiveness Reasoning Task (H3, C3, and M3)}, explain why the video is interesting in the context of the overall content. This part requires interpretive answers based on analysis, reasoning, and prior knowledge to explain why the video is counter-intuitive.

Specifically for the \textit{Magic Method Task (M4)}, we found that most annotators were not professional magicians, and even when watching instructional magic videos, it was challenging for them to provide complete and reasonable annotations for magic methods. Therefore, for M4, we only provided a small number of examples, which do not exist in the FunQA dataset. Fig. \ref{fig:M4} shows an example of M4.

We conducted strict quality reviews in real-time for the annotations, including but not limited to checking if the annotations meet the requirements, if the selected videos contain sensitive content, and if the annotation texts contain sensitive content. The annotation principles (video selection and video clip, etc) can be seen in Fig.~\ref{fig:annotation-principle}. Fig. \ref{fig:annotation-error} also illustrates some errors encountered during the annotation process and the suggested modifications provided. We assured that every annotation included in the final release of FunQA has been subjected to rigorous \textbf{multi-person, multi-round} review processes. Specifically, every piece of released FunQA version annotated content underwent scrutiny by three distinct annotators. After completing all the annotations, we conducted a consensus evaluation study to assess the objectivity of our annotations. We randomly selected 10\% of the videos and asked annotators to provide a consistency score (high consensus, medium consensus, low consensus) between their first and second annotations, considering all the previous annotations shown to them.\\
\noindent\textbf{Copyright \& License}\quad
We respect the copyright of each video. 
% For content downloaded from streaming platforms, we obtained permission from the respective uploaders. For \textit{Kasou Taisho}, we submitted a fair use request to the official website of Tokyo TV in Japan. 
Our data sharing approach is informed by the models established by previous studie like Kinetics~\cite{kay2017kinetics}, HD-VILA-100M~\cite{xue2022hdvila}, and others. Instead of providing the original raw data, we only supply the YouTube video IDs necessary for downloading the respective content.

We respect the personal identity information of everyone appearing in our videos and always strive to eliminate offensive content. By conducting strict reviews of the annotators and real-time quality checks during each annotation process, we ensure the absence of offensive content.
\subsection{Language Augmentation} \label{appendix_a: Augmentation}
For each QA pair, we first asked GPT-3.5 to translate the previous Chinese answers into English, and then let GPT-3.5 give five separate answers with the same meaning but different linguistic expressions. Afterward, we filtered out incorrect and incomplete generations. Fig. \ref{fig:prompt-trans} shows the prompt we gave to GPT-3.5.

However, as mentioned in the main text in the limitation, the translation of GPT-3.5 will be ambiguous when the original text uses Chinese harmonics and some special Chinese meanings, and this problem mostly occurs in the Title Task, we screened such problems by hand and modified them, as shown in Fig. \ref{fig:trans-err}.

\subsection{Data Statistics}\label{appendix_a: Dataset statics}
For a complete comparison between FunQA and other existing benchmarks, please see Table \ref{T:whole_main}.
We collected raw videos from multiple video platforms, including short videos, long videos, and video clips, and the detailed data can be seen in Table \ref{T:raw_statics}. The statics of FunQA dataset before GPT-3.5 extension can be seen in Table \ref{T:statics_before_gpt}. The word cloud of the all annotation word is shown in Fig. \ref{fig:wc_funqa} (a-d).

% \begin{figure}[htbp]
%   \centering
  
%   \includegraphics[width=0.22\textwidth]{figures/fig_word_cloud_funqa.png}
%   \caption{\textbf{Word cloud of FunQA free-text answer.}}
%   \label{fig:wc_funqa}
% \end{figure}


% static table----------------------------
\begin{table}[!t]
\caption{\textbf{Statistics of the FunQA raw data.} 
\label{T:raw_statics}}
\centering
\resizebox{1.0\linewidth}{!}{

\begin{tabular}{cccccc}
\hline
FunQA Subset & Type & Source & \# Videos & Avg.len (s) & Total.len (h) \\ 
\hline
\multirow{3}{*}{Humor} & \multirow{2}{*}{Daily Life (Human)} & \multirow{3}{*}{Youtube}                 & 351    & 182  & 15.8   \\
                       &                                     &                    & 1296   & 14      & 5.18    \\
                       & Nature (Animal)                     &                  & 230    & 133  & 8.52    \\ \hline
Creative             & Performance                         & Youtube            & 26     & 6060  & 43.77   \\\hline
\multirow{2}{*}{MagicQA} & Close up Magic                      & \multirow{2}{*}{Youtube} & 765    & 96  & 20.40   \\ 
                       & Camera Magic                        &                          & 334    & 152  & 14.17   \\ \hline
FunQA                    & \textit{-}                          & -                        & 3002   & 129   & 107.87  \\ \hline\\
\end{tabular}   

}
\end{table}
%static table ----------------------------------
\begin{table}[h]
\caption{\textbf{Statistics of the FunQA before GPT-3.5 extension.}}
\label{T:statics_before_gpt}
\centering
\resizebox{0.8\linewidth}{!}{
\begin{tabular}{ccccc}
\hline
Datasets & Avg length (s) & \# Clips (K) & \# QA pairs (K) & \# QA per clip\\
\hline
HumorQA & 7 & 1.8 & 7.2 & 4.0\\
CreativeQA & 48 & 0.9 & 4.5 & 5.0\\
MagicQA & 10 & 1.6 & 4.8 & 3.0\\
\hline
FunQA & 19 & 4.3 & 16.5 & 3.8\\
\hline
\end{tabular}
}
\label{tab:dataset-info}
\end{table}

\begin{figure*}
    \centering
    \begin{subfigure}[b]{0.2\textwidth}
        \centering
        \includegraphics[width=0.8\textwidth]{figures/fig_word_cloud_h.png}
        \caption{\textbf{HumorQA.}}
        \label{fig:wc_human}
    \end{subfigure}
    \begin{subfigure}[b]{0.2\textwidth}
        \centering
        \includegraphics[width=1\textwidth]{figures/fig_word_cloud_c.png}
        \caption{\textbf{CreativeQA.}}
        \label{fig:wc_creat}
    \end{subfigure}
    \begin{subfigure}[b]{0.20\textwidth}
        \centering
        \includegraphics[width=1\textwidth]{figures/fig_word_cloud_m.png}
        \caption{\textbf{MagicQA.}}
        \label{fig:wc_magic}
    \end{subfigure}
    \begin{subfigure}[b]{0.20\textwidth}
        \centering
        \includegraphics[width=1\textwidth]{figures/fig_word_cloud_funqa.png}
        \caption{\textbf{FunQA.}}
        \label{fig:wc_funqa}
    \end{subfigure}

    \caption{\textbf{Word cloud of FunQA free-text answer.}}
    \label{fig:wc_funqa}
\end{figure*}